% !TEX root = Appendix_2.tex
% =========================================================================
% APPENDIX 2 -- MICROSTRUCTURE AND COMPETING HAZARDS IN COLOMBIA
% Econometrica / Econometric Society supplement format.
% Detailed presentation of Section~2.3 results. Replication: Appendix_2/
% =========================================================================
\documentclass[ecta,nameyear,final,supplement]{econsocart}

\usepackage{bm}
\usepackage{booktabs}
\usepackage{array}
\usepackage{tabularx}
\usepackage{ragged2e}
\usepackage{longtable}
\RequirePackage[colorlinks,citecolor=blue,linkcolor=blue,urlcolor=blue,pagebackref]{hyperref}
\urlstyle{same}
\newcolumntype{Y}{>{\raggedright\arraybackslash}X}

\startlocaldefs
\numberwithin{equation}{section}
\numberwithin{table}{section}

\newcommand{\E}{\mathbb{E}}

\renewcommand{\thesection}{S\arabic{section}}
% Compact AMS-style labels: S{sec}.{sub}.{subsub}.{para}; 0 if level unused
\makeatletter
\newcommand{\@appzero}[1]{\ifnum\value{#1}=0 0\else\arabic{#1}\fi}
\newcommand{\@appsec}{S\arabic{section}}
\newcommand{\@appsub}{\@appzero{subsection}}
\newcommand{\@appsubsub}{\@appzero{subsubsection}}
\renewcommand{\thesubsection}{\@appsec.\@appsub.0.0}
\renewcommand{\thesubsubsection}{\@appsec.\@appsub.\@appsubsub.0}
\renewcommand{\theparagraph}{\@appsec.\@appsub.\@appsubsub.\arabic{paragraph}}
\setcounter{secnumdepth}{4}
\@addtoreset{subsection}{section}
\@addtoreset{subsubsection}{subsection}
\@addtoreset{paragraph}{section}
\makeatother

% APA 7-style table notes (Publication Manual, \S7.21)
\newcommand{\tabnote}{\par\addvspace{0.4em}\noindent\emph{Note.}}
\newcommand{\tabsource}{\ \textit{Source:} Authors' own calculations based on
CNMH/SIEVCAC microdata (Section~\ref{sec:data}). Software: Stata~19
(StataCorp, 2025).}
\endlocaldefs

\begin{document}

\begin{frontmatter}

\title{Supplement to ``Identifying Rational Types in Unknown Environments''\\
Appendix 2: Microstructure and Competing Hazards in Colombia}
\runtitle{Microstructure Supplement}

\begin{aug}
\author[add1,add2]{\fnms{Humberto}~\snm{Bernal}}
\address[add1]{%
\orgdiv{Economics Program},
\orgname{Universidad Colegio Mayor de Cundinamarca}}
\address[add2]{%
\orgdiv{Ph.D. in Economics},
\orgname{Universidad de los Andes}}
\ead[label=e1]{hbernal@uniandes.edu.co}
\ead[label=e2]{hbernalc@universidadmayor.edu.co}
\end{aug}

\begin{funding}
Economist, professor of the Economics Program at the Universidad Colegio
Mayor de Cundinamarca and graduate of the Ph.D.\ in Economics from the
Universidad de los Andes, Colombia.
\end{funding}

\end{frontmatter}

\noindent
This supplement presents in full detail the microeconometric results
summarized in Section~2.3 of the main paper (``Microstructure and competing
hazards in Colombia''). It documents the CNMH/SIEVCAC data, the construction
of the analytical sample, the sample-selection test, the multinomial logit
(MNL) model of kidnapping outcomes, specification diagnostics and robustness
checks, and the cause-specific Cox proportional hazards models of resolution
timing. All estimates were reproduced from the Stata scripts in the
\texttt{Appendix\_2/} replication package (Section~\ref{sec:repro}).

% ============================================================================
\section{Data and sample construction}
\label{sec:data}
% ============================================================================

Individual kidnapping records come from the Information System on Events of
Violence in the Colombian Armed Conflict (SIEVCAC), maintained by the
National Center of Historical Memory (CNMH) \citep{cnmh_datos,CNMH2013}.
The replication package includes the merged case--victim file
\path{stata/data/Data_merge.dta} ($N=39{,}369$ victim-level records) and the
estimation-ready checkpoint \path{stata/data/analytical_sample.dta} ($n=1{,}125$
after outcome filters, variable construction, and dropping eight cases with
missing geographic zone).

\subsection{Universe and sequential filters}

Table~\ref{tab:sample-flow} summarizes the sample flow reported in the main
text. The universe comprises all documented kidnapping victims in the merged
CNMH export. The first filter retains extortionate kidnappings
(\texttt{TipodeSecuestro = EXTORSIVO}) with a valid event date. The second
filter requires an observable liberation outcome (excluding \texttt{ND},
\texttt{OTRO}, and \texttt{NA} in \texttt{TipodeLiberación}). The third
filter requires complete covariates for estimation; eight cases with missing
geographic zone are dropped, yielding $n=1{,}125$.

\begin{table}[t]
\centering
\caption{Analytical Sample Construction}
\label{tab:sample-flow}
\begin{tabular}{@{}lrr@{}}
\toprule
Stage & $n$ & Share of prior stage \\
\midrule
Universe (merged CNMH records) & 39,369 & --- \\
Valid event date & 36,916 & 93.8\% \\
Extortionate typology & 11,270 & 30.5\% \\
Observable outcome & 1,133 & 10.1\% \\
Complete covariates & 1,125 & 99.3\% \\
\bottomrule
\end{tabular}
\tabnote\ Valid event date: non-missing \texttt{Fecha} and \texttt{Año}$\neq 0$
after parsing \path{Data_merge.dta}. Extortionate share of universe:
$11{,}270/39{,}369=28.6\%$. Estimation share: $1{,}125/39{,}369=2.9\%$. Eight
cases with missing geographic zone are dropped at the final stage.
\tabsource
\end{table}

\subsection{Outcome variable}

The dependent variable \texttt{Y\_Resultado} classifies each extortionate
episode into four mutually exclusive outcomes (reference category: Payment):
\begin{itemize}
\item \textbf{Payment} ($Y=1$): ransom or exchange paid
  (\texttt{PAGO}, \texttt{CANJE}, and related codes).
\item \textbf{Rescue} ($Y=3$): liberation by state forces (\texttt{RESCATE}).
\item \textbf{Escape or release without payment} ($Y=0$):
  \texttt{FUGA} or release under external pressure without ransom.
\item \textbf{Death} ($Y=2$): victim killed in captivity
  (\texttt{MUERTO EN CAUTIVERIO}, \texttt{ASESINADO}).
\end{itemize}
Table~\ref{tab:outcome-dist} reports the distribution in the estimation
sample.

\begin{table}[t]
\centering
\caption{Outcome Distribution in the Estimation Sample}
\label{tab:outcome-dist}
\begin{tabular}{@{}lrr@{}}
\toprule
Outcome & Count & Share \\
\midrule
Payment (reference) & 512 & 45.5\% \\
Rescue & 471 & 41.9\% \\
Escape or release without payment & 121 & 10.8\% \\
Death in captivity & 21 & 1.9\% \\
\midrule
Total & 1,125 & 100.0\% \\
\bottomrule
\end{tabular}
\tabnote\ Estimation sample $n=1{,}125$. Payment is the MNL reference category.
\tabsource
\end{table}

\subsection{Covariates}

The covariate vector $\mathbf{X}_{i,m}\in\mathbb{R}^{26}$ in the main text
comprises four blocks (all coded as indicator variables with stated
reference categories):
\begin{enumerate}
\item \textbf{Perpetrator} ($\mathbf{G}_i$): FARC (ref.), ELN, paramilitaries,
  common delinquency/other.
\item \textbf{Victim profile} ($\mathbf{V}_i$): high income (ref.), middle
  class, vulnerable, public sector, unknown.
\item \textbf{Context} ($\mathbf{C}_{i,m}$): geographic zone (metropolis ref.,
  Andean, Caribbean, Pacific/conflict, East/jungle) and historical period
  (pre-1998 ref., 1998--2002, 2003--2010, post-2011).
\item \textbf{Tactics} ($\mathbf{T}_i$): modus operandi, collective vs.\
  individual structure, GAULA intervention, victim sex, captivity duration
  category (express $\le 2$ days ref., short $\le 30$ days, long $>30$ days,
  missing), and $\ln(\text{total victims per case})$.
\end{enumerate}
Standard errors are clustered at the municipal level (373~clusters). Variable
construction rules are implemented in \path{microstructure_replication.do}
and \path{sample_selection_robustness.do}.

% ============================================================================
\section{Sample-selection test}
\label{sec:selection}
% ============================================================================

Because only 10.1\% of extortionate cases have observable outcomes, a binary
logistic model was estimated on the extortionate stratum with complete
attrition covariates ($n=11{,}253$ of $11{,}270$ extortionate cases with valid
date), with dependent variable \texttt{Missing\_Outcome}$=1$ when the liberation
type is unobserved \citep{LittleRubin2020}. Table~\ref{tab:attrition} summarizes
the main findings as odds ratios with cluster-robust standard errors by
municipality.

\begin{table}[t]
\centering
\caption{Attrition Logistic Regression Results}
\label{tab:attrition}
\small
\begin{tabular}{@{}lccr@{}}
\toprule
Covariate & OR & $p$-value & Significance \\
\midrule
Non-metropolitan zones & 1.9--2.2 & $<0.001$ & 1\% \\
Period 1998--2002 & 1.76 & $0.001$ & 1\% \\
Victim profile & --- & $>0.10$ & --- \\
$\ln$(total victims per case) & --- & $>0.10$ & --- \\
\bottomrule
\end{tabular}
\tabnote\ Estimated on extortionate cases with complete attrition covariates
($n=11{,}253$; $11{,}270$ extortionate cases with valid date). Stata drops 17
observations with missing covariates. Cluster-robust standard errors by
municipality. OR $>1$ indicates higher odds of an unobservable outcome.
Non-metropolitan zones comprise Andean, Caribbean, Pacific, and East/jungle
categories. Victim-profile and scale covariates are not predictive at the 10\%
level; geography and the 1998--2002 period show systematic attrition. The
Significance column reports the highest conventional level (1\%, 5\%, 10\%);
``---'' denotes not significant at 10\%.
\tabsource
\end{table}

The test rejects complete randomness of missing outcomes while supporting
the internal validity of the structural estimators: substantive covariates do
not predict missingness, and the remaining bias is territorial (metropolitan
overrepresentation).

% ============================================================================
\section{Multinomial logit model}
\label{sec:mnl}
% ============================================================================

\subsection{Specification}

With Payment as the base category ($Y=1$), the MNL estimates
\begin{equation}
\ln\!\left(\frac{\Pr(Y_{i,m}=j \mid \mathbf{X}_{i,m})}{\Pr(Y_{i,m}=1 \mid \mathbf{X}_{i,m})}\right)
= \alpha_j + \mathbf{X}_{i,m}'\boldsymbol{\beta}_j,
\quad j\in\{0,2,3\}.
\label{eq:mnl-app}
\end{equation}
The main specification uses cluster-robust standard errors at the municipality
level \citep{Wooldridge2010}.

\subsection{Model fit}

The model is jointly significant: Wald $\chi^2(78)=5{,}250.34$, $p<0.001$;
McFadden pseudo-$R^2=0.18$.

\subsection{Full estimation results}

Table~\ref{tab:mnl-full} reports the complete MNL specification: all 26
covariates for the three non-payment outcomes (Escape, Death, Rescue), with
cluster-robust standard errors, $z$-statistics, $p$-values, and 95\%
confidence intervals. The baseline outcome is Payment ($Y=1$).

% Full multinomial logit table (Stata mlogit; N=1,125). Matches Section 2.3 Spanish text.
{\footnotesize
\renewcommand{\arraystretch}{1.15}
\begin{longtable}{@{}l r r r r c c@{}}
\caption{Multinomial Logistic Regression Results (Full Model)}
\label{tab:mnl-full} \\
\toprule
\textbf{Outcome / Predictors} & \textbf{RRR} & \textbf{Std.\ Err.} & \textbf{$z$} & \textbf{$P>|z|$} & \multicolumn{2}{c}{\textbf{[95\% Conf.\ Interval]}} \\
\midrule
\endfirsthead
\multicolumn{7}{c}{\textbf{Table~\thetable{} --- continued from previous page}} \\
\toprule
\textbf{Outcome / Predictors} & \textbf{RRR} & \textbf{Std.\ Err.} & \textbf{$z$} & \textbf{$P>|z|$} & \multicolumn{2}{c}{\textbf{[95\% Conf.\ Interval]}} \\
\midrule
\endhead
\midrule
\multicolumn{7}{r}{\textit{Continued on next page}} \\
\bottomrule
\endfoot
\bottomrule
\endlastfoot

\multicolumn{7}{l}{\textbf{Outcome: Escape or release without payment}} \\
\textit{Perpetrator group (ref.\ FARC)} & & & & & & \\
ELN & 0.800 & 0.284 & $-0.63$ & 0.531 & 0.399 & 1.606 \\
Paramilitaries & 0.932 & 0.419 & $-0.16$ & 0.876 & 0.386 & 2.252 \\
Common delinquency/other & 0.493 & 0.233 & $-1.50$ & 0.135 & 0.195 & 1.246 \\
\addlinespace
\textit{Victim profile (ref.\ upper class)} & & & & & & \\
Middle class & 1.868 & 0.749 & 1.56 & 0.119 & 0.852 & 4.097 \\
Vulnerable population & 0.624 & 0.278 & $-1.06$ & 0.291 & 0.260 & 1.496 \\
Public sector & 3.917 & 1.842 & 2.90 & 0.004 & 1.559 & 9.844 \\
Unknown & 1.469 & 0.519 & 1.09 & 0.276 & 0.736 & 2.935 \\
\addlinespace
\textit{Geographic zone (ref.\ metropolis)} & & & & & & \\
Andean & 0.637 & 0.284 & $-1.01$ & 0.313 & 0.266 & 1.528 \\
Caribbean & 0.496 & 0.238 & $-1.46$ & 0.143 & 0.194 & 1.268 \\
Pacific/conflict zone & 0.554 & 0.269 & $-1.22$ & 0.224 & 0.214 & 1.436 \\
East/jungle & 0.727 & 0.399 & $-0.58$ & 0.561 & 0.248 & 2.133 \\
\addlinespace
\textit{Modus operandi (ref.\ direct attack)} & & & & & & \\
Illegal roadblock & 2.287 & 1.187 & 1.59 & 0.111 & 0.827 & 6.323 \\
Military incursion & 1.274 & 0.627 & 0.49 & 0.622 & 0.486 & 3.341 \\
Deception & 0.878 & 0.483 & $-0.24$ & 0.814 & 0.299 & 2.583 \\
Unknown & 1.444 & 0.438 & 1.21 & 0.226 & 0.797 & 2.615 \\
\addlinespace
\textit{Kidnapping structure} & & & & & & \\
Group structure (vs.\ individual) & 0.760 & 0.308 & $-0.68$ & 0.498 & 0.343 & 1.683 \\
\addlinespace
\textit{GAULA intervention} & & & & & & \\
GAULA intervention (vs.\ none) & 0.942 & 0.350 & $-0.16$ & 0.873 & 0.455 & 1.952 \\
\addlinespace
\textit{Victim sex (ref.\ male)} & & & & & & \\
Female & 0.625 & 0.260 & $-1.13$ & 0.259 & 0.276 & 1.414 \\
No information & 2.038 & 3.405 & 0.43 & 0.670 & 0.077 & 53.832 \\
\addlinespace
\textit{Duration (ref.\ express $\le 2$ days)} & & & & & & \\
Short ($\le 30$ days) & 0.322 & 0.198 & $-1.85$ & 0.065 & 0.096 & 1.072 \\
Long ($>30$ days) & 0.235 & 0.156 & $-2.18$ & 0.029 & 0.064 & 0.861 \\
Missing duration & 0.464 & 0.282 & $-1.26$ & 0.207 & 0.141 & 1.529 \\
\addlinespace
$\ln$(total victims per case) & 0.764 & 0.197 & $-1.04$ & 0.297 & 0.461 & 1.267 \\
\addlinespace
\textit{Historical period (ref.\ pre-1998)} & & & & & & \\
Crisis/Cagu\'an (1998--2002) & 1.506 & 0.675 & 0.91 & 0.360 & 0.626 & 3.624 \\
Democratic security (2003--2010) & 1.235 & 0.587 & 0.44 & 0.658 & 0.486 & 3.135 \\
Peace process/current (2011+) & 11.283 & 7.227 & 3.78 & 0.000 & 3.215 & 39.596 \\
\addlinespace
Constant & 0.486 & 0.401 & $-0.87$ & 0.382 & 0.097 & 2.447 \\
\midrule

\multicolumn{7}{l}{\textbf{Outcome: Death in captivity}} \\
\textit{Perpetrator group (ref.\ FARC)} & & & & & & \\
ELN & 4.500 & 6.741 & 1.00 & 0.315 & 0.239 & 84.787 \\
Paramilitaries & 20.306 & 25.564 & 2.39 & 0.017 & 1.722 & 239.473 \\
Common delinquency/other & 63.629 & 85.557 & 3.09 & 0.002 & 4.562 & 887.547 \\
\addlinespace
\textit{Victim profile (ref.\ upper class)} & & & & & & \\
Middle class & 0.515 & 0.876 & $-0.39$ & 0.696 & 0.018 & 14.455 \\
Vulnerable population & 7.056 & 5.195 & 2.65 & 0.008 & 1.667 & 29.873 \\
Public sector & 0.000 & 0.000 & $-13.09$ & 0.000 & 0.000 & 0.000 \\
Unknown & 1.219 & 0.998 & 0.24 & 0.809 & 0.245 & 6.062 \\
\addlinespace
\textit{Geographic zone (ref.\ metropolis)} & & & & & & \\
Andean & 0.560 & 0.604 & $-0.54$ & 0.591 & 0.067 & 4.643 \\
Caribbean & 0.000 & 0.000 & $-15.61$ & 0.000 & 0.000 & 0.000 \\
Pacific/conflict zone & 0.456 & 0.521 & $-0.69$ & 0.492 & 0.049 & 4.284 \\
East/jungle & 0.624 & 0.733 & $-0.40$ & 0.688 & 0.062 & 6.246 \\
\addlinespace
\textit{Modus operandi (ref.\ direct attack)} & & & & & & \\
Illegal roadblock & 0.000 & 0.000 & $-16.21$ & 0.000 & 0.000 & 0.000 \\
Military incursion & 0.000 & 0.000 & $-22.26$ & 0.000 & 0.000 & 0.000 \\
Deception & 0.537 & 0.698 & $-0.48$ & 0.632 & 0.042 & 6.863 \\
Unknown & 1.244 & 0.736 & 0.37 & 0.712 & 0.390 & 3.967 \\
\addlinespace
\textit{Kidnapping structure} & & & & & & \\
Group structure (vs.\ individual) & 11.744 & 8.468 & 3.42 & 0.001 & 2.858 & 48.262 \\
\addlinespace
\textit{GAULA intervention} & & & & & & \\
GAULA intervention (vs.\ none) & 0.536 & 0.512 & $-0.65$ & 0.514 & 0.082 & 3.489 \\
\addlinespace
\textit{Victim sex (ref.\ male)} & & & & & & \\
Female & 2.450 & 1.831 & 1.20 & 0.230 & 0.566 & 10.600 \\
No information & 47.771 & 83.378 & 2.22 & 0.027 & 1.561 & 1461.557 \\
\addlinespace
\textit{Duration (ref.\ express $\le 2$ days)} & & & & & & \\
Short ($\le 30$ days) & 0.580 & 0.546 & $-0.58$ & 0.562 & 0.092 & 3.666 \\
Long ($>30$ days) & 0.477 & 0.545 & $-0.65$ & 0.517 & 0.051 & 4.474 \\
Missing duration & 0.361 & 0.420 & $-0.88$ & 0.381 & 0.037 & 3.522 \\
\addlinespace
$\ln$(total victims per case) & 0.084 & 0.054 & $-3.90$ & 0.000 & 0.024 & 0.293 \\
\addlinespace
\textit{Historical period (ref.\ pre-1998)} & & & & & & \\
Crisis/Cagu\'an (1998--2002) & 1.230 & 1.215 & 0.21 & 0.834 & 0.177 & 8.529 \\
Democratic security (2003--2010) & 1.405 & 1.397 & 0.34 & 0.732 & 0.200 & 9.856 \\
Peace process/current (2011+) & 10.008 & 9.700 & 2.38 & 0.017 & 1.497 & 66.890 \\
\addlinespace
Constant & 0.004 & 0.010 & $-1.98$ & 0.047 & 0.000 & 0.932 \\
\midrule

\multicolumn{7}{l}{\textbf{Outcome: Rescue}} \\
\textit{Perpetrator group (ref.\ FARC)} & & & & & & \\
ELN & 0.840 & 0.198 & $-0.74$ & 0.459 & 0.529 & 1.334 \\
Paramilitaries & 0.680 & 0.250 & $-1.05$ & 0.295 & 0.330 & 1.399 \\
Common delinquency/other & 0.494 & 0.174 & $-2.00$ & 0.045 & 0.247 & 0.984 \\
\addlinespace
\textit{Victim profile (ref.\ upper class)} & & & & & & \\
Middle class & 1.277 & 0.334 & 0.93 & 0.350 & 0.765 & 2.130 \\
Vulnerable population & 0.800 & 0.221 & $-0.81$ & 0.420 & 0.465 & 1.376 \\
Public sector & 0.654 & 0.281 & $-0.99$ & 0.323 & 0.282 & 1.518 \\
Unknown & 0.749 & 0.197 & $-1.10$ & 0.272 & 0.447 & 1.255 \\
\addlinespace
\textit{Geographic zone (ref.\ metropolis)} & & & & & & \\
Andean & 0.800 & 0.329 & $-0.54$ & 0.587 & 0.357 & 1.792 \\
Caribbean & 0.426 & 0.194 & $-1.87$ & 0.061 & 0.174 & 1.041 \\
Pacific/conflict zone & 0.656 & 0.311 & $-0.89$ & 0.374 & 0.259 & 1.661 \\
East/jungle & 0.740 & 0.332 & $-0.67$ & 0.501 & 0.307 & 1.782 \\
\addlinespace
\textit{Modus operandi (ref.\ direct attack)} & & & & & & \\
Illegal roadblock & 1.513 & 0.667 & 0.94 & 0.348 & 0.638 & 3.589 \\
Military incursion & 1.112 & 0.388 & 0.31 & 0.760 & 0.561 & 2.205 \\
Deception & 1.232 & 0.559 & 0.46 & 0.645 & 0.507 & 2.998 \\
Unknown & 2.077 & 0.518 & 2.93 & 0.003 & 1.274 & 3.386 \\
\addlinespace
\textit{Kidnapping structure} & & & & & & \\
Group structure (vs.\ individual) & 0.972 & 0.280 & $-0.10$ & 0.921 & 0.552 & 1.710 \\
\addlinespace
\textit{GAULA intervention} & & & & & & \\
GAULA intervention (vs.\ none) & 1.390 & 0.481 & 0.95 & 0.341 & 0.706 & 2.737 \\
\addlinespace
\textit{Victim sex (ref.\ male)} & & & & & & \\
Female & 1.894 & 0.415 & 2.92 & 0.004 & 1.233 & 2.910 \\
No information & 7.430 & 10.051 & 1.48 & 0.138 & 0.524 & 105.297 \\
\addlinespace
\textit{Duration (ref.\ express $\le 2$ days)} & & & & & & \\
Short ($\le 30$ days) & 0.183 & 0.085 & $-3.65$ & 0.000 & 0.074 & 0.455 \\
Long ($>30$ days) & 0.079 & 0.043 & $-4.69$ & 0.000 & 0.027 & 0.228 \\
Missing duration & 0.448 & 0.171 & $-2.11$ & 0.035 & 0.213 & 0.945 \\
\addlinespace
$\ln$(total victims per case) & 1.116 & 0.197 & 0.62 & 0.534 & 0.789 & 1.578 \\
\addlinespace
\textit{Historical period (ref.\ pre-1998)} & & & & & & \\
Crisis/Cagu\'an (1998--2002) & 1.028 & 0.345 & 0.08 & 0.935 & 0.532 & 1.986 \\
Democratic security (2003--2010) & 1.554 & 0.595 & 1.15 & 0.249 & 0.734 & 3.290 \\
Peace process/current (2011+) & 3.849 & 1.906 & 2.72 & 0.006 & 1.459 & 10.157 \\
\addlinespace
Constant & 2.314 & 1.429 & 1.36 & 0.174 & 0.690 & 7.765 \\
\midrule
\addlinespace
\textit{$N$} & \multicolumn{6}{c}{1,125} \\
\textit{Wald $\chi^2(78)$} & \multicolumn{6}{c}{5,250.34} \\
\textit{Prob $>\chi^2$} & \multicolumn{6}{c}{0.0000} \\
\textit{Pseudo $R^2$} & \multicolumn{6}{c}{0.1847} \\
\end{longtable}
}
\tabnote\ Relative risk ratios (RRR) reported. Baseline outcome is Payment.
Cluster-robust standard errors at the municipality level (373 clusters;
$n=1{,}125$). $p$-values in parentheses are two-sided. Significance at 5\%:
$|z|>1.96$.
\tabsource


\subsection{Key relative risk ratios}

Table~\ref{tab:mnl-key} summarizes the three headline findings in Section~2.3
of the main paper, expressed as relative risk ratios (RRR $=\exp(\hat\beta)$).
The corresponding rows in Table~\ref{tab:mnl-full} are Peace process/current
(Escape), public-sector victim (Escape), and short/long duration (Rescue).

\begin{table}[t]
\centering
\caption{Selected Multinomial Logit Relative Risk Ratios}
\label{tab:mnl-key}
\small
\begin{tabular}{@{}llccr@{}}
\toprule
Outcome & Covariate & RRR & $p$-value & Significance \\
\midrule
Escape & Post-2011 period & 11.28 & $<0.01$ & 1\% \\
Escape & Public-sector victim & 3.92 & $<0.01$ & 1\% \\
Rescue & Short captivity ($\le 30$~d) & 0.18 & $<0.001$ & 1\% \\
Rescue & Long captivity ($>30$~d) & 0.08 & $<0.001$ & 1\% \\
\bottomrule
\end{tabular}
\tabnote\ Base category: Payment. RRR $=\exp(\hat\beta)$. Duration
categories are relative to express captivities ($\le 2$~days). Period
contrast: post-2011 vs.\ pre-1998. The Significance column reports the
highest conventional level (1\%, 5\%, 10\%); ``---'' denotes not significant
at 10\%.
\tabsource
\end{table}

\subsection{Predicted probabilities (margins)}

Table~\ref{tab:mnl-margins} reports average marginal predicted probabilities
from the main MNL (unconditional margins).

\begin{table}[t]
\centering
\caption{Average Predicted Probabilities from the Main MNL}
\label{tab:mnl-margins}
\small
\begin{tabular}{@{}llr@{}}
\toprule
Dimension & Category & $\Pr(\cdot)$ \\
\midrule
\multicolumn{3}{@{}l@{}}{\textit{Rescue by geographic zone}} \\
& Metropolis & 0.465 \\
& Andean & 0.447 \\
& Pacific & 0.414 \\
& East/jungle & 0.424 \\
& Caribbean & 0.343 \\
\midrule
\multicolumn{3}{@{}l@{}}{\textit{Death by perpetrator group}} \\
& FARC & 0.002 \\
& ELN & 0.007 \\
& Paramilitaries & 0.026 \\
& Common delinquency & 0.071 \\
\bottomrule
\end{tabular}
\tabnote\ Unconditional margins from \texttt{margins} after the main MNL
($n=1{,}125$). Rescue probabilities decline away from metropolitan nodes;
death probabilities rise as logistical retention capacity falls.
\tabsource
\end{table}

% ============================================================================
\section{MNL diagnostics and robustness}
\label{sec:mnl-diag}
% ============================================================================

\subsection{IIA test}

The Hausman--McFadden \citep{HausmanMcFadden1984} SUEST generalization
\citep{Weesie1999} compares the full four-outcome model with a restricted
model that excludes Death. The test does not reject independence of
irrelevant alternatives: $\chi^2(1)=0.24$, $p=0.625$.

\subsection{Coefficient stability}

Three specifications are compared: (i)~main model with municipal clusters;
(ii)~robust standard errors without clustering; and (iii)~specification
without duration categories.

Table~\ref{tab:robust} reports two coefficients from the Death equation.

\begin{table}[t]
\centering
\caption{MNL Robustness Coefficients}
\label{tab:robust}
\small
\begin{tabular}{@{}lccc@{}}
\toprule
 & \multicolumn{3}{c}{$\hat\beta$ (Death equation)} \\
\cmidrule(lr){2-4}
Covariate & Main & Robust SE & No duration \\
\midrule
Common delinquency (vs.\ FARC) & 4.15 & 4.15 & 4.15 \\
Collective structure (vs.\ individual) & 2.46 & 2.46 & 2.46 \\
\bottomrule
\end{tabular}
\tabnote\ Municipal clusters in Main and No duration; heteroskedasticity-
robust standard errors in Robust SE. Coefficients are stable across
specifications: both remain positive and significant at the 1\% level in all
three columns ($p<0.01$; No duration $p<0.001$).
\tabsource
\end{table}

\subsection{Death-rate heterogeneity}

Raw death rates in the estimation sample are 8.0\% for common delinquency
and 0.2\% for the FARC. A two-sample proportion test (FARC vs.\ common
delinquency) rejects equal rates: $z=-5.71$, $p<0.001$.

% ============================================================================
\section{Cause-specific Cox proportional hazards}
\label{sec:cox}
% ============================================================================

\subsection{Specification}

Following \citet{Cox1972} and the competing-risks framework of
\citet{Bradburn2003}, four cause-specific Cox models are estimated on
\texttt{DíasdeCautiverio} (days in captivity). In each model, one outcome
defines the event of interest and the remaining three outcomes are treated
as right-censored. The cause-specific hazard is
\begin{equation}
h_j(t \mid \mathbf{X}) = h_{0j}(t)\exp\!\left(\mathbf{X}'\tilde{\boldsymbol{\beta}}_j\right),
\qquad j\in\{\text{Payment, Rescue, Death, Escape}\},
\label{eq:cox-app}
\end{equation}
with the same covariate blocks as the MNL but excluding duration categories,
which would be collinear with the time axis. The Cox subsample comprises
$N=461$ cases with positive captivity duration and complete covariates (602
cases with missing duration and 64 zero-day captivities are excluded; Breslow
ties). FARC is the reference perpetrator group. Cluster-robust standard
errors are reported.

\subsection{Estimation results}

Table~\ref{tab:cox-full} reports all four cause-specific equations (Payment,
Rescue, Death, and Escape/release without payment). Table~\ref{tab:cox-hr}
lists the coefficients emphasized in Section~2.3 of the main paper.

% Cause-specific Cox: four equations (N=461). Stata 19, cox_only.do.
{\footnotesize
\renewcommand{\arraystretch}{1.05}
\begin{longtable}{lcccc}
\caption{Cause-Specific Cox Proportional Hazards Results (Four Equations)}
\label{tab:cox-full} \\
\toprule
& \textbf{Payment} & \textbf{Rescue} & \textbf{Death} & \textbf{Escape} \\
\textbf{Predictors} & (HR) & (HR) & (HR) & (HR) \\
\midrule
\endfirsthead
\multicolumn{5}{c}{\textbf{Table~\thetable{} --- continued from previous page}} \\
\toprule
& \textbf{Payment} & \textbf{Rescue} & \textbf{Death} & \textbf{Escape} \\
\textbf{Predictors} & (HR) & (HR) & (HR) & (HR) \\
\midrule
\endhead
\midrule
\multicolumn{5}{r}{\textit{Continued on next page}} \\
\bottomrule
\endfoot
\bottomrule
\endlastfoot

\multicolumn{5}{l}{\textbf{\textit{Perpetrator group (ref.\ FARC)}}} \\
ELN & 1.186 & 1.057 & 3.039 & 0.569 \\
    & {\scriptsize (0.317)} & {\scriptsize (0.837)} & {\scriptsize (0.404)} & {\scriptsize (0.225)} \\
Paramilitaries & 4.284*** & 1.262 & 12.211* & 7.547*** \\
               & {\scriptsize (0.000)} & {\scriptsize (0.530)} & {\scriptsize (0.049)} & {\scriptsize (0.000)} \\
Common delinquency & 1.713* & 0.935 & 49.988** & 1.157 \\
                   & {\scriptsize (0.021)} & {\scriptsize (0.862)} & {\scriptsize (0.004)} & {\scriptsize (0.812)} \\
\midrule

\multicolumn{5}{l}{\textbf{\textit{Victim profile (ref.\ upper class)}}} \\
Middle class & 0.780 & 1.071 & 0.000*** & 0.680 \\
             & {\scriptsize (0.274)} & {\scriptsize (0.823)} & {\scriptsize (0.000)} & {\scriptsize (0.467)} \\
Vulnerable population & 1.542* & 0.512 & 13.108* & 0.676 \\
                      & {\scriptsize (0.046)} & {\scriptsize (0.053)} & {\scriptsize (0.030)} & {\scriptsize (0.512)} \\
Public sector & 0.199*** & 0.237*** & 0.000*** & 2.133 \\
              & {\scriptsize (0.000)} & {\scriptsize (0.001)} & {\scriptsize (0.000)} & {\scriptsize (0.218)} \\
\midrule

\multicolumn{5}{l}{\textbf{\textit{Geographic zone (ref.\ metropolis)}}} \\
Andean & 1.290 & 1.709 & 0.126 & 0.445 \\
       & {\scriptsize (0.450)} & {\scriptsize (0.465)} & {\scriptsize (0.052)} & {\scriptsize (0.209)} \\
Caribbean & 1.170 & 1.024 & 0.000*** & 0.523 \\
          & {\scriptsize (0.649)} & {\scriptsize (0.975)} & {\scriptsize (0.000)} & {\scriptsize (0.304)} \\
East/jungle & 0.815 & 0.625 & 0.231 & 0.674 \\
            & {\scriptsize (0.644)} & {\scriptsize (0.530)} & {\scriptsize (0.117)} & {\scriptsize (0.579)} \\
\midrule

\multicolumn{5}{l}{\textbf{\textit{Modus operandi (ref.\ direct attack)}}} \\
Illegal roadblock & 1.780* & 2.339 & 0.000*** & 5.044** \\
                  & {\scriptsize (0.012)} & {\scriptsize (0.069)} & {\scriptsize (0.000)} & {\scriptsize (0.003)} \\
Deception & 1.330 & 3.275*** & 0.000*** & 0.555 \\
          & {\scriptsize (0.309)} & {\scriptsize (0.000)} & {\scriptsize (0.000)} & {\scriptsize (0.421)} \\
\midrule

\multicolumn{5}{l}{\textbf{\textit{Kidnapping characteristics}}} \\
Group structure (vs.\ individual) & 0.950 & 0.577 & 3.736 & 3.226** \\
                                  & {\scriptsize (0.801)} & {\scriptsize (0.114)} & {\scriptsize (0.331)} & {\scriptsize (0.008)} \\
GAULA intervention (vs.\ none) & 0.716 & 0.597 & 1.172 & 0.976 \\
                               & {\scriptsize (0.267)} & {\scriptsize (0.244)} & {\scriptsize (0.932)} & {\scriptsize (0.961)} \\
$\ln$(total victims per case) & 1.205 & 1.593* & 0.659 & 0.295*** \\
                              & {\scriptsize (0.191)} & {\scriptsize (0.052)} & {\scriptsize (0.739)} & {\scriptsize (0.000)} \\
\midrule

\multicolumn{5}{l}{\textbf{\textit{Historical period (ref.\ pre-1998)}}} \\
Crisis/Cagu\'an (1998--2002) & 1.112 & 0.677 & 1.366 & 0.401 \\
                             & {\scriptsize (0.514)} & {\scriptsize (0.243)} & {\scriptsize (0.878)} & {\scriptsize (0.150)} \\
Democratic security (2003--2010) & 0.735 & 2.022* & 9.186 & 0.220 \\
                                 & {\scriptsize (0.125)} & {\scriptsize (0.036)} & {\scriptsize (0.108)} & {\scriptsize (0.063)} \\
Peace process/current (2011+) & 0.510* & 4.863*** & 35.559** & 4.886** \\
                              & {\scriptsize (0.035)} & {\scriptsize (0.000)} & {\scriptsize (0.008)} & {\scriptsize (0.002)} \\
\midrule

\multicolumn{5}{l}{\textit{Model statistics}} \\
Observations ($N$) & 461 & 461 & 461 & 461 \\
Events & 277 & 120 & 13 & 51 \\
Wald $\chi^2$ & 117.57 & 182.68 & --- & 116.45 \\
Prob $>\chi^2$ & 0.0000 & 0.0000 & 0.0000 & 0.0000 \\
\end{longtable}
}
\tabnote\ Hazard ratios (HR) with cluster-robust $p$-values in parentheses
(two-sided). * $p<.05$; ** $p<.01$; *** $p<.001$. Duration categories are
excluded (collinear with the time axis). The Cox subsample ($N=461$) drops
602 cases with missing duration and 64 zero-day captivities (Breslow ties).
Payment and Rescue coefficients match Section~2.3 of the main paper; Death
and Escape equations use the same covariate vector (\path{cox_only.do}). The
Death equation has only 13 events: several HRs hit numerical bounds (reported
as 0.000) owing to quasi-complete separation; its Wald statistic is omitted
from interpretation. Pacific zone, unknown victim profile, military
incursion, unknown modus, and sex indicators are estimated but omitted for
parity with the main-text table layout.
\tabsource


\begin{table}[t]
\centering
\caption{Selected Cause-Specific Cox Hazard Ratios}
\label{tab:cox-hr}
\small
\begin{tabular}{@{}llccr@{}}
\toprule
Equation & Covariate & HR & $p$-value & Significance \\
\midrule
\multicolumn{5}{@{}l@{}}{\textit{Payment equation}} \\
& Paramilitaries & 4.28 & $<0.001$ & 1\% \\
& Common delinquency & 1.71 & $0.021$ & 5\% \\
& ELN & 1.19 & $0.317$ & --- \\
\midrule
\multicolumn{5}{@{}l@{}}{\textit{Payment and Rescue equations}} \\
& Public-sector victim (Payment) & 0.20 & $<0.001$ & 1\% \\
& Public-sector victim (Rescue) & 0.24 & $0.001$ & 1\% \\
\midrule
\multicolumn{5}{@{}l@{}}{\textit{Post-2011 period}} \\
& Post-2011 (Payment) & 0.51 & $0.035$ & 5\% \\
& Post-2011 (Rescue) & 4.86 & $<0.001$ & 1\% \\
\bottomrule
\end{tabular}
\tabnote\ Rows match Table~\ref{tab:cox-full} (rounded for the narrative).
Reference perpetrator group: FARC. HR $=\exp(\tilde\beta)$; HR $>1$ indicates
faster resolution, HR $<1$ slower resolution. The Significance column reports
the highest conventional level (1\%, 5\%, 10\%); ``---'' denotes not
significant at 10\%. Payment speed falls 49\% post-2011 ($1-0.51$); Rescue
hazard multiplies by ${\approx}4.86$.
\tabsource
\end{table}

\subsection{GAULA intervention and cumulative rescue}

Unconditional margins from the Rescue Cox model show that GAULA intervention
raises the predicted rescue probability at each duration category
(Table~\ref{tab:cox-gaula}).

Cumulative rescue probabilities by day~100 appear in
Table~\ref{tab:cox-cumul} (exported in \path{outputs/Datos_Graficas_Cox.xlsx}).

\begin{table}[t]
\centering
\caption{Rescue Probability by GAULA Intervention and Duration}
\label{tab:cox-gaula}
\small
\begin{tabular}{@{}lccc@{}}
\toprule
GAULA intervention & Express & Short & Long \\
\midrule
No & 0.64 & 0.28 & 0.14 \\
Yes & 0.72 & 0.35 & 0.19 \\
\bottomrule
\end{tabular}
\tabnote\ Unconditional margins from the Rescue cause-specific Cox model.
Duration categories: express ($\le 2$~days), short ($\le 30$~days), long
($>30$~days).
\tabsource
\end{table}

\begin{table}[t]
\centering
\caption{Cumulative Rescue Probability at Day 100}
\label{tab:cox-cumul}
\small
\begin{tabular}{@{}lr@{}}
\toprule
Perpetrator group & $\Pr(\text{Rescue by day 100})$ \\
\midrule
Paramilitaries & 0.33 \\
ELN & 0.29 \\
FARC & 0.27 \\
\bottomrule
\end{tabular}
\tabnote\ Computed from cause-specific survival curves ($1-S_t$) at day
$t=98$ (maximum grid point in \path{stcurve} with \texttt{range(0 100)}) in
\path{outputs/cumul_rescue_d100.dta}, exported to
\path{outputs/Datos_Graficas_Cox.xlsx} by \path{microstructure_replication.do}.
\tabsource
\end{table}

\subsection{Mortality}

The MNL yields unconditional average predicted death probabilities of 0.002
(FARC) and 0.071 (common delinquency) in Table~\ref{tab:mnl-margins}. The
corresponding cause-specific Cox margins reported in Section~2.3 of the main
paper are 0.002 and 0.079 (${\approx}35\times$ and ${\approx}40\times$,
respectively). Both specifications preserve the same ordering: mortality rises
as logistical retention capacity falls.

% ============================================================================
\section{Relevance for the mechanism}
\label{sec:relevance}
% ============================================================================

\paragraph{Scope.}
This section maps the microeconometric results to the mechanism-design model
(Section~2.3 of the main paper, ``Microstructure and competing hazards in
Colombia''). It introduces no new estimates. Every empirical quantity cited
below is reported in Sections~S1--S7 and can be verified from the replication
package using \path{README.txt} at the package root.

\begin{table}[t]
\centering
\caption{Mechanism Inputs and Replication Files}
\label{tab:mechanism-map}
\small
\begin{tabularx}{\textwidth}{@{}>{\raggedright\arraybackslash}Y
  >{\raggedright\arraybackslash}Y
  >{\raggedright\arraybackslash}Y@{}}
\toprule
Mechanism object & Main-text result & Package file \\
\midrule
Prior $\mu_0$ over types $\theta_K$ & MNL outcome distribution &
  \path{microstructure_replication.do} \\
Terminal payment technology & RRR by victim profile &
  \path{sample_selection_robustness.do} \\
Arrival rates $\lambda_j(\theta_K,t)$ & Cox cause-specific HRs &
  \path{microstructure_replication.do} \\
GAULA effectiveness $\gamma_t$ & GAULA $\times$ duration margins &
  \path{outputs/Datos_Graficas_Cox.xlsx} \\
Sample validity & Attrition logit, IIA, robustness &
  \path{sample_selection_robustness.do} \\
\bottomrule
\end{tabularx}
\tabnote\ The MNL supplies the static type map; the Cox models supply
dynamic resolution hazards that feed the daily public signal $\zeta_t$ in the
MDG process. All paths are relative to \path{Appendix_2/stata/}.
\tabsource
\end{table}

% ============================================================================
\section{Reproducibility}
\label{sec:repro}
% ============================================================================

\begin{sloppypar}
All results in this supplement are reproduced by Stata scripts in the
\texttt{Appendix\_2/} replication package. Set the working directory to
\texttt{Appendix\_2/} and run \path{replicate.do}. See \path{README.txt} at the
package root for step-by-step instructions. Requirements: Stata~19 or later.
Estimated runtime: 5--15 minutes. Batch mode:
\texttt{stata -b -q do replicate.do}.
\end{sloppypar}

\begin{enumerate}
\item \path{replicate.do} runs the full pipeline in \path{stata/}:
  \path{microstructure_replication.do}, \path{sample_selection_robustness.do},
  and \path{cox_only.do}.
\item Individual scripts live in \path{stata/}; each includes
  \path{_setup_paths.do} and locates \path{stata/data/Data_merge.dta}
  automatically.
\item Generated outputs: \path{stata/outputs/analytical_sample.dta},
  \path{stata/outputs/cox_appendix_rows.txt},
  \path{stata/outputs/cox_cumul_day100.txt}, and log files.
\item Cox curves: \path{stata/outputs/Datos_Graficas_Cox.xlsx} and
  \path{stata/outputs/cumul_rescue_d100.dta} (Tables~\ref{tab:cox-gaula}
  and \ref{tab:cox-cumul}); regenerated by \path{microstructure_replication.do}.
\item (Optional) Compile \path{tex/Appendix_2.pdf} from \path{tex/} with
  \texttt{pdflatex} and \texttt{bibtex}.
\end{enumerate}

\begin{sloppypar}
Section~\ref{sec:relevance} (Table~\ref{tab:mechanism-map}) maps
microeconometric outputs to mechanism inputs. Raw CNMH exports used to build
\path{stata/data/Data_merge.dta} are documented in \path{README.txt}.
\end{sloppypar}

\bibliographystyle{econsoc}
\bibliography{references}

\end{document}
